\documentclass{article}
\usepackage{amsmath}

\title{Unbiased estimation via randomization - Application to
Option Pricing}
\author{Alic Kaufmann}
\date{31.05.2016}

\begin{document}

\maketitle
(Q1)
1.
\begin{align*}
    \|Y_l -Y\|_2^2
    &= E[Y_l-Y]^2\\
    &= E[Y_l^2 - 2Y_l Y + Y^2]\\
    &= E[Y_l^2]-2E[Y_l]E[Y]+E[Y^2]\quad\text{ by independence of }Y_l\text{ and }Y\\
    &\to E[Y^2]-2 E[Y]^2 + E[Y^2]\\
    &= 2\underbrace{(E[Y^2]-E[Y]^2)}_{Var[Y]}\\
    &=2 Var[Y] \not\to 0\quad\text{ if }Var[Y]\neq0
\end{align*}

By assumption, $N$ is almost surely finite, i.e. $$P(N=\infty)=0.$$

But $$\{N\ge 1\}\supset \{N\ge 2\}\supset\{N\ge3\}\supset\ldots$$
and  $\cap_{l\ge 1}\{N\ge l\}=\{N=\infty\}.$ so by continuity of probability,
$$0=P(N=\infty)=\lim_{l\to\infty}P(N\ge l), $$ so for $l$ big enough
$$\frac{\|Y_{l-1}-Y\|_2^2}{P(N\ge l)}\ge \frac{Var[Y]}{P(N\ge l)}.$$
But since $P(N\ge l)\to0$ we have $$\frac{Var[Y]}{P(N\ge l)}\to\infty$$ so
$$\sum_{l=1}^{\infty} \frac{\|Y_{l-1}-Y_l\|_2^2}{P(N\ge l)}=\infty$$
2. We want to show

\begin{equation}\label{dummy}
\frac{\|Y_{l-1}-Y\|_2^2}{P(N\ge l)}<\infty.
\end{equation}
By the strong order $p$ assumption and the fact that we have a geometric time step, we have
\begin{align*}
\frac{\|Y_{l-1}-Y\|_2^2}{P(N\ge l)}&\le \frac{\left(C\left(\frac{T}{2^{l-1}}\right)^p\right)^2}{q^l} \\
&= \underbrace{CT^{2p}2^{2p}}_{\tilde{C}}\left(\frac{1}{2^{2p}q}\right)^l.
\end{align*}
In order for this geometric series to converge, we need $\frac{1}{2^{2p}q}<1$, i.e. $q>2^{-2p}.$

We now procede to show that $\bar{Z}$ has finite variance.
By definition,  $Var[\bar{Z}]=E[\bar{Z^2}]-\left[E\bar{Z}\right]^2.$ But by Jensen's inequaliy, we have
$$E[\bar{Z}]=E[Y]\le (E[Y^2])^{1/2}<\infty\quad\text{ because }Y\in L^2(\Omega).$$
So we just need to show $E[\bar{Z}^2]<\infty.$
But 
\begin{align*}
    E[\bar{Z}^2]
    &=\sum_{l=0}^{\infty} \frac{\bar{\nu}_l}{P(N\ge l)}\\
    &=\sum_{l=0}^{\infty} \frac{\|Y_{l-1}-Y\|_2^2-\|Y_l-Y\|_2^2}{P(N\ge l)}\\
    &\le\sum_{l=0}^{\infty} \frac{\|Y_{l-1}-Y\|_2^2}{P(N\ge l)}<\infty\quad\text{ by }\ref{dummy}
\end{align*}

(Q2)
1.
\begin{align*}
    E[\bar{\tau}]\\
    &= E\left[\sum_{j=1}^N \bar{t}_j\right]\\
    &= E\left[\sum_{j=1}^{\infty}\bar{t}_j1(N\ge j)\right]\\
    &= \sum_{j=1}^{\infty}\bar{t}_jE\left[1(N\ge j)\right]\quad\text{ by Tonelli}\\
    &= \sum_{j=1}^{\infty} \bar{t}_jP(N\ge j)
\end{align*}
2. Since $g(\alpha F) = g(F)$,
$$\min_F \left(\sum_{j=0}^{\infty}\frac{\beta_j}{F_j}\right)\left(\sum_{j=0}^{\infty}t_j F_j\right)$$
is equivalent to solving
\begin{align*}
    &\min_F &\sum_{j=0}^{\infty}\frac{\beta_j}{F_j}\\
    &\text{subject to } &\sum_{j=0}^{\infty}t_jF_j=C
\end{align*}

This is a constrained optimization problem so we use lagrange multipliers:
\begin{align*}
\mathcal{L}(F,\lambda)=\sum_{j=0}^{\infty}\frac{\beta_j}{F_j}+\lambda\left(\sum_{j=0}^{\infty}t_j F_j-C\right)\\
\end{align*}

\begin{align*}
\frac{\partial\mathcal{L}}{\partial F_i} &= 0\\
\frac{-\beta_i}{F_i^2}+\lambda t_i&=0\\
F_i &= \frac{1}{\sqrt{\lambda}}\sqrt{\frac{\beta_i}{t_i}}
\end{align*}
In order that $F_0=1$, we need $$\frac{1}{\sqrt{\lambda}}\sqrt{\frac{\beta_0}{t_0}}=1$$
so $\lambda = \beta_0/t_0$, i.e. $$F_i=\frac{\sqrt{\beta_i/t_i}}{\sqrt{\beta_0/t_0}}=:F^*_i.$$
So if $F^*$ is feasible (as stated in the assumption of the quesiton) we have that
$F^*$ satisfies the KKT conditions. Moreover, the problem is convex so $F^*$ is an optimal solution.




3. \begin{align*}
\|\Delta_j\|_2
&=\|Y_j-Y_{j-1}\|\\
&\le \underbrace{\|Y_j - Y\|_2}_{C2^{-pj}} + \underbrace{\|Y-Y_{j-1}\|_2}_{C2^{-p(j-1)}}\\
&\le (C + C2^p)2^{-pj}
\end{align*}
so
\begin{align*}
E\left[\Delta_j^2\right]&=\|\Delta_j\|_2^2\\
&=\underbrace{\left(C + C2^p\right)^2}_{\tilde{C}} 2^{-2pj}\\
&=O(2^{-2pj})
\end{align*}
4.In order to get the best balance between the variance and the cost, we note that we need
to consider the optimal $F^*_j=P(N\ge j)$.

We know $F^*_j \sim \sqrt{\frac{\nu_j}{t_j}} \sim \sqrt{\frac{2^{-2pj}}{2^j}}=2^{-j(p+\frac12)}$.
Since we assumed that the survival distribution is of the form $P(N\ge l)=2^{-rl}$,
we have $r^*=p + \frac12$ is balancing the variance and the cost.

Lets the conditions for $E[\bar{\tau}]$ and $Var[\bar{Z}]$ to be finite.
\begin{align*}
    E[\bar{\tau}]
    &= \sum_{j=0}^{\infty} t_j P(N\ge j)\\
    &\le \sum_{j=0}^{\infty}c2^j2^{-rj}\\
    &= \sum_{j=0}^{\infty}C 2^{j(1-r)}
\end{align*}
In order for the series to converge, we need $j(1-r)<0$, i.e. $r>1$.

By the same argument as in (Q1.2), $Var[\bar{Z}]<\infty$ is equivalent to $E\left[\bar{Z}^2\right]<\infty$.
But 
\begin{align*}
    E[\bar{Z}^2]
    &= \sum_{l=0}^{\infty}\frac{\bar{\nu}_l}{P(N\ge l)}\\
    &= \sum_{l=0}^{\infty}\frac{C2^{-2pj}}{2^{-rj}}\\
    &= \sum_{l=0}^{\infty}C 2^{j(r-2p)}
\end{align*}
and in order for the series to converge, we need $j(r-2p)<0$, i.e. $r<2p$.
\end{document}
